\section{Designing a finite-label release}

The cellwise criterion in \cref{obj:thm:universal-point-identification} makes label design an allocation problem. A release with \leanref{sym:K}{\ensuremath{K}} labels groups score values, and \cref{obj:def:worst-case-ambiguity} measures the resulting worst-case ATE ambiguity. A lower bound on score density supports a comparison across score laws.

\begin{assumptionv}{ass:density-lower}[Lower bound on score density]
The score law \(H\) has density \(f\) on the score interval \(\mathcal E\), with
\[
H(de)=f(e)\,de,\qquad f(e)\ge m_f
\]
for Lebesgue-almost every \(e\in\mathcal E\).
\end{assumptionv}

The density \leanref{sym:f}{\ensuremath{f}} and its floor \leanref{sym:mf}{\ensuremath{m_f}}, when the floor is positive, give the standard compact-support positive-density condition used in quantization \citep{LemaireMontesPages2020}. A positive floor assigns every subinterval a baseline amount of score mass, which permits a uniform lower bound on ambiguity.

The release budget, its optimal ambiguity, and the common comparison class are defined as follows.

\begin{definitionv}{def:k-label-releases}[$K$-label release class $\mathcal G_K$]
\[
\mathcal G_K
=
\bigl\{g:\mathcal E\to\{1,\ldots,K\}:g\text{ is measurable}\bigr\}.
\]
\end{definitionv}

The release family \leanref{sym:GK}{\ensuremath{\mathcal G_K}} allows any measurable partition with at most \(K\) nonempty label cells. Its design criterion takes the infimum of worst-case ambiguity over that family.

\begin{definitionv}{def:optimal-k-ambiguity}[Optimal $K$-label ambiguity $\Delta_K(H)$]
\[
\Delta_K(H)=\inf_{g\in\mathcal G_K}D_H(g).
\]
\end{definitionv}

The optimal ambiguity \leanref{sym:DeltaK}{\ensuremath{\Delta_K(H)}} expresses the best worst-case ATE resolution available at the given label budget. For a uniform comparison, the density floor defines a class of score laws on the common overlap interval.

\begin{definitionv}{def:score-lower-density-class}[Lower-density score class $\mathcal H^{\mathrm{lb}}_{\varepsilon,m_f}$]
Under \cref{obj:ass:density-lower}, the class of score laws is
\[
\mathcal H^{\mathrm{lb}}_{\varepsilon,m_f}
=
\left\{
H\in\mathcal P(\mathcal E):
H(de)=f(e)\,de,\quad
f(e)\ge m_f\ \text{for Lebesgue-almost every }e\in\mathcal E
\right\}.
\]
\end{definitionv}

The class is nonempty only in the feasible range \(0<m_f\le (1-2\varepsilon)^{-1}\); this is the range in which the uniform lower bound below has content. Within this class, the label budget yields matching inverse-budget bounds. The constructive upper bound applies to every score law on the overlap interval.

\begin{theoremv}{thm:k-label-rate}[Bounds for optimal label ambiguity]
Suppose:
\begin{itemize}
\item \textbf{(Overlap.)} The condition in \cref{obj:ass:overlap} holds for \(\varepsilon\).
\item \textbf{(Density floor.)} \(m_f>0\).
\item \textbf{(Label budget.)} \(K\geq 1\) is an integer.
\end{itemize}
Let \(\ell=1-2\varepsilon\) and \(\mathcal E=[\varepsilon,1-\varepsilon]\). For every \(H\in\mathcal H^{\mathrm{lb}}_{\varepsilon,m_f}\) as defined in \cref{obj:def:score-lower-density-class},
\[
\frac{m_f\ell^2}{2(1-\varepsilon)K}\leq\Delta_K(H).
\]
For every probability law \(H\) on \(\mathcal E\), the equal-width release \(g:\mathcal E\to\mathcal R_K=\{1,\ldots,K\}\) defined by
\[
g(e)=1+\min\left\{\left\lfloor\frac{K(e-\varepsilon)}{\ell}\right\rfloor,K-1\right\}
\]
is measurable and satisfies
\[
\Delta_K(H)\leq\frac{\ell}{2\varepsilon(1-\varepsilon)K},
\qquad
D_H(g)\leq\frac{\ell}{2\varepsilon(1-\varepsilon)K}.
\]
\end{theoremv}

A positive density floor gives an inverse-budget lower bound on ambiguity for every \(K\)-label design. Equal-width cells give a matching upper bound for the optimal ambiguity under that floor, and an inverse-budget upper bound for any score law, including laws with atoms. The displayed constants thus translate a label budget into a worst-case ATE-width guarantee before outcomes are analyzed.

For a continuous, strictly positive score density, the leading coefficient has an exact law-specific form. An interval-cell release attains that coefficient asymptotically.

\begin{theoremv}{thm:sharp-high-resolution-constant}[Sharp high resolution constant]
Let \(\mathcal E=[\varepsilon,1-\varepsilon]\) be the score interval, and let \(H\) be a probability law on \(\mathcal E\). Suppose:
\begin{itemize}
\item \textbf{(Overlap.)} The overlap condition in \cref{obj:ass:overlap} holds.
\item \textbf{(Density.)} \(H(de)=f(e)\,de\) on \(\mathcal E\).
\item \textbf{(Regularity.)} The density \(f\) is continuous and strictly positive on \(\mathcal E\).
\end{itemize}
Let \(\Delta_K(H)\) and \(D_H(g)\) be the optimal \(K\)-label ambiguity and the worst-case ambiguity defined in \cref{obj:def:optimal-k-ambiguity,obj:def:worst-case-ambiguity}. For every positive integer \(K\), there is a measurable release \(g_K:\mathcal E\to\{0,\ldots,K-1\}\) whose label cells are intervals, such that
\[
\lim_{K\to\infty}K\Delta_K(H)
=
\lim_{K\to\infty}K D_H(g_K)
=
\frac14\left(
\int_{\varepsilon}^{1-\varepsilon}
\sqrt{\frac{f(e)}{e(1-e)}}\,de
\right)^2.
\]
\end{theoremv}

The theorem writes the attaining labels as \(0,\ldots,K-1\), whereas \(\mathcal G_K\) uses \(1,\ldots,K\). The bijection \(r\mapsto r+1\) preserves every label cell and hence preserves \(D_H(g_K)\), so the displayed construction is a member of the release class after relabeling.

The integrand weights score regions by their density and by the reciprocal-assignment factor \(1/[e(1-e)]\). It suggests an ordered interval allocation with boundaries at approximately equal increments of the cumulative weight \(\int_{\varepsilon}^{e}\sqrt{f(t)/[t(1-t)]}\,dt\): regions with greater weight receive narrower cells and more labels. The constant in \cref{obj:thm:sharp-high-resolution-constant} quantifies the leading ambiguity under optimal allocation, and its interval-cell sequence attains that value.
